%% kapa: exam=1 round=30 year=2019 period=1 form=A subject=civil_law subject_ko=민법 seq=1 number=1
%% source: Q-Net 감정평가사 기출문제 (cst00309) artlSeq=5205205 file=extracted/1008/5205205/2019년 제30회 감정평가사 1차시험 문제지/2019년 제30회 감정평가사 1차시험 문제지/2019년도 제30회 감정평가사 1차 1교시 A형.hwp
%% conversion: kordoc 4.24.51 -> tools/gen1.py (auto); review=proofread-sonnet; flags=none
\begin{question}{1}
민법의 법원(法源)에 관한 설명으로 옳은 것은? (다툼이 있으면 판례에 따름)
\begin{choices}
  \item 관습법에 앞서 적용되는 법률이란 국회에서 제정된 법률만을 말한다.
  \item 관습법에 의한 분묘기지권은 더 이상 인정되지 않는다.
  \item 판례는 관습법과 사실인 관습을 구별하지 않는다.
  \item 상급법원 재판에서의 판단은 해당 사건에 관하여 하급심을 기속한다.
  \item 헌법재판소의 결정은 그것이 민사에 관한 것이라도 민법의 법원으로 되지 않는다.
\end{choices}
\end{question}
\answer{4}
